% Generated draft: requires mathematical and source review.
\documentclass[11pt]{article}
\usepackage[margin=25mm]{geometry}
\usepackage{fontspec}
\setmainfont{Latin Modern Roman}
\usepackage{amsmath,amssymb,mathtools,unicode-math}
\setmathfont{Latin Modern Math}
\usepackage{enumitem,needspace,xurl,hyperref}
\setlength{\parindent}{0pt}
\setlength{\parskip}{5pt}
\setlength{\emergencystretch}{3em}
\newcommand{\N}{\mathbb{N}}
\newcommand{\Z}{\mathbb{Z}}
\newcommand{\Q}{\mathbb{Q}}
\newcommand{\R}{\mathbb{R}}
\newcommand{\C}{\mathbb{C}}
\newcommand{\F}{\mathbb{F}}
\newcommand{\A}{\mathbb{A}}
\newcommand{\PP}{\mathbb{P}}
\newcommand{\E}{\mathbb{E}}
\newcommand{\eps}{\varepsilon}
\newcommand{\dd}{\,\mathrm{d}}
\begin{document}
\section*{UnsolvedMath ID: 800816}
\textbf{Source problem code:} catalogue entry
\subsection{Definitions and mathematical statement}
Let \(M\) be a compact Riemannian \(3\)-manifold with Riemannian volume measure \(\mu\). Write \(\mathscr D^s_\mu(M)\) for the group of volume-preserving diffeomorphisms of Sobolev class \(H^s\), and let \(e\) denote its identity. In the usual Sobolev manifold framework one takes \(s>5/2\); this regularity convention is explanatory and is not specified in the supplied source excerpt. The excerpt also does not specify whether \(M\) has boundary or which boundary conditions apply. The construction below is understood subject to a choice of these conventions for which the hydrodynamic exponential map is defined.

The right-invariant \(L^2\) metric is determined at the identity by
\[
\langle u,v\rangle_{L^2}=\int_M g_x(u(x),v(x))\,\mathrm d\mu(x),
\qquad u,v\in T_e\mathscr D^s_\mu(M).
\]
Let \(\mathcal U\subset T_e\mathscr D^s_\mu(M)\) be the open domain of initial velocities whose geodesics exist through time \(1\), and define
\[
\operatorname{Exp}^{L^2}_e:\mathcal U\longrightarrow\mathscr D^s_\mu(M),
\qquad \operatorname{Exp}^{L^2}_e(u)=\eta_u(1),
\]
where \(\eta_u(0)=e\) and \(\dot\eta_u(0)=u\). Here \(L^2\) specifies the metric; the differential is considered between tangent spaces with their \(H^s\) topology. For \(u\in\mathcal U\), put
\[
A_u=D\operatorname{Exp}^{L^2}_e(u):T_e\mathscr D^s_\mu(M)\longrightarrow T_{\operatorname{Exp}^{L^2}_e(u)}\mathscr D^s_\mu(M).
\]
A bounded linear operator is Fredholm if its kernel is finite-dimensional, its range is closed, and its cokernel is finite-dimensional. The exponential map is Fredholm on its domain if every \(A_u\) is Fredholm.

\textbf{Problem (supplied source, Problem 15).} Quantify the failure of the Fredholm property of \(\operatorname{Exp}^{L^2}_e\) for a compact Riemannian \(3\)-manifold \(M\). In terms of the preceding definitions, the requested subject is the non-Fredholm behavior of the operators \(A_u\). The excerpt does not specify a numerical measure of failure, a locus of initial velocities, or a uniformity requirement in \(M\) or \(u\); consequently it does not determine a unique quantified assertion. Kernel dimension, closedness of the range, and cokernel dimension describe the Fredholm criteria, but the source does not prescribe particular bounds or classifications for them.

The supplied motivation is a qualitative comparison: the structure of singularities of the hydrodynamic exponential map in two dimensions resembles that of a smooth map between finite-dimensional manifolds. No precise two-dimensional theorem or hypotheses for that comparison are supplied.
\subsection{Short English statement}
For the \(L^2\) geodesic exponential map on the group of volume-preserving Sobolev diffeomorphisms of a compact Riemannian \(3\)-manifold, describe quantitatively how its differential fails to be Fredholm. The supplied question leaves the intended quantification, Sobolev regularity, and boundary conventions unspecified. The registry labels the problem open; this is a supplied status, not an independently verified assessment.
\subsection{Sources}
\begin{itemize}
\item Supplied local UnsolvedMath registry record: permanent ID \(800816\), source problem code \texttt{catalogue entry}, title \emph{Geometric Hydrodynamics in Open Problems --- Problem 15}. The record attributes the problem to Khesin, Misiolek, and Modin, with proposed year \(2022\), and labels its status \emph{open}.
\item Khesin, Misiolek, and Modin, \emph{Geometric Hydrodynamics in Open Problems} (2022), Problem 15. Supplied source URL: \url{https://arxiv.org/abs/2205.01143}. The record identifies the extraction method as \texttt{arxiv\_latex} and gives the access date \texttt{2026-10-16}. This access date and provenance are supplied metadata; the source was not independently visited or verified for this draft.
\item Supplied upstream revision: \texttt{d371ee917ac9647fbb06cd6fc529deec6f3dda82}. Registry creation and update timestamps are both \texttt{2026-10-16T14:00:00Z}.
\end{itemize}
\end{document}
